% Section 2 - std::thread
% Roberto Masocco <roberto.masocco@uniroma2.it>
% April 22, 2026

% ### std::thread ###
\section{std::thread}
\graphicspath{{figs/section2/}}

% --- std::thread ---
\begin{frame}{std::thread}{From POSIX to C++}
	\justifying
	You already know threads from your OS course: a thread is an \textbg{independent flow of execution} within a process, sharing its \textbg{address space} and \textbg{resources}.\\
	\bigskip
	In C/POSIX, threads are managed through the \texttt{pthread} API: \texttt{pthread\_create}, \texttt{pthread\_join}, \texttt{pthread\_detach} — functional but verbose.\\
	\bigskip
	C++11 introduced \texttt{std::thread}: a \textbg{portable}, \textbg{RAII-aware wrapper} around the platform's threading primitives.\\
	The underlying model is identical:
	\begin{itemize}
		\item threads share the process memory space;
		\item each thread has its own stack;
		\item scheduling is managed by the operating system kernel.
	\end{itemize}
	\begin{block}{}
		\centering
		\textbf{Same model, cleaner interface, fewer mistakes.}
	\end{block}
\end{frame}
\begin{frame}[fragile]{std::thread}{Creating threads}
	\justifying
	An \texttt{std::thread} is constructed by passing any \textbg{callable} and its arguments.\\
	The thread starts executing \textbg{immediately} upon construction, arguments are \textbg{copied} by default.
	\begin{columns}
		\column{.9\textwidth}
		\begin{lstlisting}[style=codebeamer, language=C++, caption=Creating threads from a function and a lambda.]
#include <iostream>
#include <thread>

void greet(const std::string & name) {
  std::cout << "Hello from " << name << "\n";
}

std::thread t1(greet, "thread_1");

std::thread t2([]() {std::cout << "Hello from lambda\n";});
    \end{lstlisting}
	\end{columns}
\end{frame}
\begin{frame}[fragile]{std::thread}{\texttt{join()} and \texttt{detach()}}
	\justifying
	Every \texttt{std::thread} must be either \textbg{joined} or \textbg{detached} before it is destroyed.\\
	\bigskip
	\texttt{t.join()} blocks the calling thread until \texttt{t} finishes — the standard way to \textbg{wait for completion}.\\
	\bigskip
	\texttt{t.detach()} releases ownership: the thread runs independently and its resources are reclaimed by the OS when it finishes — use sparingly and only for \textbg{fire-and-forget} tasks.
	\begin{columns}
		\column{.9\textwidth}
		\begin{lstlisting}[style=codebeamer, language=C++, caption=Joining threads.]
std::thread t1(greet, "worker_1");
std::thread t2(greet, "worker_2");

t1.detach(); // do not wait for t1 to finish
t2.join();   // wait for t2 to finish
    \end{lstlisting}
	\end{columns}
\end{frame}
